\subsection{Data processing}\label{subsec:mm-data-processing}

In this case study, we utilised the DongNiao International Birds 10000 Dataset (DIB-10K)
by \textcite{Mei2020Dongniao}, which comprises over 4.8 million images representing
10,922 bird species, following the IOC 10.1 taxonomy \parencite{Gill2021IOC}. This
extensive dataset encompasses a wide variety of bird species, morphological variants,
postures, and gestures \parencite{Mei2020Dongniao}. Despite their efforts in manual
review and correction, the dataset contains numerous duplicate and erroneous images,
necessitating a comprehensive data cleaning process.

For the data cleaning, we employed the pHash method to detect near-duplicate images
by comparing these hashes \parencite{Farid2021Overview}. Duplicates spanning multiple
categories are completely removed, and for intracategory duplicates, only one copy is
retained in its category.

To address erroneous images with non-avian subjects, we utilised the pretrained Faster
Region-based Convolutional Neural Network (Faster R-CNN,~\cite{Ren2017Faster}) in
the torchvision library \parencite{Ansel2024Pytorch}. For images where no birds were
detected, we implemented a two-tiered approach: images from high-volume categories
($\geq200$ images) were automatically deleted, while those from low-volume categories
($<200$ images) were moved to a separate directory for manual inspection.
After a manual thorough review, valid images were reintegrated into the dataset.

\subsection{Model and its basic characteristics}\label{subsec:r-model}

The task of recognising over 10,000 bird species is a
fine-grained visual categorisation (FGVC) problem, where the goal is
to identify minor differences between highly similar categories. To
tackle this challenge, we adopted the MetaFGNet model proposed by \textcite{Zhang2018Fine},
which is specifically designed for FGVC tasks. We utilised their
pretrained weights "LBird-31\_checkpoint" as the foundation for
transfer learning on the processed DIB-10K dataset. In the training process,
all species are regarded as equal categories, without any \textit{a priori} taxonimic
knowledge being introduced.

For the training process, we employed a server equipped with an RTX 4090D
GPU, 90GB of memory, and a 15-core
Intel\textsuperscript{\textregistered} Xeon\textsuperscript{\textregistered}
Platinum 8474C processor from the \href{https://www.autodl.com/docs/}{AutoDL platform}.
This high-performance setup allowed for efficient processing and model training.
The training procedure was conducted for 32 epochs.

We reassigned the orders and families of all species (categories) to align with the
IOC 15.1 \parencite{Gill2025IOC}. The weights are extracted from the final fully
connected layer (fc) of the ResNet34 model. Weights of each species were regarded as a
512-dimensional vector representing the morphological traits. They are reduced to the
lowest boundary that can explain 80\% of all variances for subsequent analyses.
Gradient-weighted Class Activation Mapping (Grad-CAM,~\cite{Selvaraju2020Grad}) is
employed to visualise the model's attention on images.

The model achieved 90.8\% accuracy on the training set and 87.6\% accuracy
on the validation set. Grad-CAM reveals that the network's attention is focused
on birds, effectively ignoring complex backgrounds and occlusions. The attention
maps cover the entire bird, indicating that the extracted phenotypic vectors
can be regarded as representations of the shape, plumage, and colour of species (\cref{fig:gradcam}).

\begin{figure}[htbp]
    \centering
    \includegraphics[width=\textwidth]{gradcam}
    \caption{Grad-CAM reveals that the network's attention is consistently
    focused on birds, ignoring backgrounds and occlusions.}
    \label{fig:gradcam}
\end{figure}

The global mean pairwise angle across all birds is approximately
1.57066 radians ($\approx \frac{\pi}{2}$), with a low variance
($\approx 0.00968$). In contrast, the intra-taxa distributions (families and orders) exhibit
significantly lower mean angles ($\approx1.2$) but higher variances
(peaking at $\approx0.02\textminus0.03$) (\cref{fig:angles}).

A strong positive correlation exists between taxa size
and mean pairwise angle (Orders: $\rho=0.91$; Families: $\rho=0.75$).
Mean pairwise angle shows weak correlation with angle variance
(Orders: $\rho=0.33$; Families: $\rho=0.47$).
Conversely, taxa size shows weak or no correlation with angle variance
(Orders: $\rho=0.24,p>0.05$; Families: $\rho=0.34$).

\begin{figure}[htbp]
    \centering
    \includegraphics[width=\textwidth]{angles_analysis}
    \caption{(a-f) The relationship between taxa size and mean pairwise angle,
        taxa size and pairwise angle variance, as well as mean pairwise angle and pairwise
        angle variance at both order and family levels. Spearman's rank coefficients ($\rho$)
        and p-values are listed for each figure; (g-h) the distribution of mean pairwise angle
        and pairwise angle variance of families and orders, the mean values of all birds are
        marked as the red dash line.}
    \label{fig:angles}
\end{figure}

\subsection{Similarity clustering}\label{subsec:r-similarity-clustering}

The clustering process results in a comprehensive hierarchical clustering output.
The agglomerative clustering technique applied to the cosine similarity
measures of the weight vectors yielded a dendrogram that illustrates the
relationships between the different avian species based on their
morphological features learned by the ResNet34 model.

In the taxonomic consistence analysis, we identified a total of 391
branches with high taxonomic consistency at the family level and 94
branches at the order level. Additionally, we found 474 family-level
outlier species and 533 order-level outlier species (\cref{fig:cluster}).

\begin{figure}[htbp]
  \centering
  \begin{minipage}{0.45\textwidth}
    \centering
    \includegraphics[width=1.05\textwidth]{cluster-1}
  \end{minipage}
  \hfill
  \begin{minipage}{0.45\textwidth}
    \centering
    \includegraphics[width=\textwidth]{cluster-2}
  \end{minipage}
  \caption{The unrooted clustering result with high-purity branches collapsed.}
\label{fig:cluster}
\end{figure}

\subsection{Disparity analysis}\label{subsec:r-disparity-analysis}

The Spearman's rank correlation was employed to examine
the relationship between diversity (species richness) and morphological
disparity (spherical variance), where taxa with $N<2$ were excluded. The Akaike
information criterion (AIC) was employed to assess and compare four models:

\begin{enumerate}
    \item Power law: $f(x) = 1 - x^b$, the coefficient of $x^b$
        is fixed as 1 to satisfy the boundary condition $f(1) = 0$,
        reflecting the theoretical expectation;
    \item Stretched exponential model: $f(x) = 1 - \exp(-\lambda \cdot (x - 1)^\beta)$;
    \item Hill equation: $f(x) = \frac{(x - 1)^n}{k + (x - 1)^n}$;
    \item Logarithmic rational model: $f(x) = \frac{\ln(x)}{k + \ln(x)}$.
\end{enumerate}

All analyses were conducted at both order and family levels. Subsequently, the residuals
of the disparity of every taxon are calculated, five taxa with the highest and five with
the lowest residuals are listed.

Additionally, we explore the relationship between taxa size and mean pairwise angle,
taxa size and pairwise angle variance, as well as mean pairwise angle and pairwise
angle variance at both order and family levels. For size vs variance analysis,
taxa with $N<3$ were excluded to avoid mathematical artefacts
($\text{variance} = 0$ for $N=1 \text{or} 2$). For the other two analysis (mean angle),
taxa with $N<2$ were excluded.

Regarding the relationship between order-level disparity and diversity, the
Spearman's coefficient is at 0.966, with a p-value of $1.45\times10^{-24}$.
For families, the Spearman's coefficient was 0.908,
with a p-value of $3.56\times10^{-83}$.

The stretched exponential model receives the lowest AIC (-251.33) at the order
level, followed by the Hill equation (-250.02), without significant difference
($\Delta AIC = 1.31$). The models are respectively $f(x) = 1 - \exp(-0.1420(x - 1)^{0.3081})$
and $f(x) = \frac{(x - 1)^{0.3973}}{7.4957 + (x - 1)^{0.3973}}$.
Orders with the highest disparity residuals are Falconiformes, Cuculiformes,
Pelecaniformes, Mesitornithiformes (highest first), and Piciformes.
Caprimulgiformes, Procellariiformes, Apterygiformes, Struthioniformes, and Psittaciformes
(lowest first, the same applies below) show the least disaprity residuals (\cref{subfig:models_order}).

\begin{figure}[htbp]
  \centering
  \begin{subfigure}{0.48\textwidth}
    \centering
    \includegraphics[width=\textwidth]{model_comparison_order}
    \caption{}
    \label{subfig:models_order}
  \end{subfigure}
  \hfill
  \begin{subfigure}{0.48\textwidth}
    \centering
    \includegraphics[width=\textwidth]{model_comparison_family}
    \caption{}
    \label{subfig:models_family}
  \end{subfigure}
  \caption{Comparison of models of the phenotypic disparity of taxa (orders or famalies) as a
  function of diversity. Showing scatter plot of all taxa, graphs of four models, names of
  the five taxa with the highest and the five with the lowest phenotypic disparity.}
  \label{fig:models_comparison}
\end{figure}

For families, the power law receives the lowest AIC (-1201.08), followed by the
Hill equation (-1199.76), without significant difference
($\Delta AIC = 1.32$). The models are respectively $f(x) = 1 - x^{-0.1429}$
and $f(x) = \frac{(x - 1)^{0.3723}}{6.0303 + (x - 1)^{0.3723}}$.
Mitrospingidae, Vangidae, Oreoicidae, Ptiliogonatidae, and Modulatricidae show the
highest disparity residuals, and Apodidae, Rhinocryptidae, Caprimulgidae, Hydrobatidae,
and Procellariidae show the lowest (\cref{subfig:models_family}).

\subsection{Disparity through time}\label{subsec:r-disparity-through-time}

The disparity-through-time (DTT) analysis reveals an early-burst pattern for the evolution of
avian visual morphospace. As shown in \cref{fig:dtt}, the empirical data of relative disparity
deviates extremely from the expectation of Brownian Motion null simulation.

Following the K-Pg mass extinction ($\sim 66$ Mya), the empirical disparity immediately increases
at an intense rate. This curve is consistently higher than the mean null expectation throughout
the Paleogene and Neogene, resulting in a positive Morphological Disparity Index (MDI).

Crucially, the empirical curve achieves approximately 50\% of the modern disparity shortly after
the K-Pg mass extinction, followed by a relative deceleration towards the present. Conversely,
the BM model requires significantly more time to reach the same disparity.

\begin{figure}[htbp]
    \centering
    \includegraphics[width=0.65\textwidth]{dtt_null_test_combined}
    \caption{Relative disparity-through-time (DTT) plot for avian visual morphospace.
    The solid blue line represents the empirical disparity, while the solid grey line
    indicates the mean expectation under a multivariate Brownian Motion (BM) null model.
    Both empirical and simulational data are normalised relative to the disparity of
    the extant time (0 Mya). The K-Pg boundary (66 Mya) is marked as the red dash line.}
    \label{fig:dtt}
\end{figure}

\subsection{Macroevolution of birds}\label{subsec:macroevolution-of-birds}

The Spearman's rank correlation demonstrates that species richness is the primary driver of morphospace expansion.
The strong positive correlation between taxon size and hyperspherical variance (Orders: $\rho = 0.966$; Families: $\rho = 0.908$)
confirms that as clades diversify, they tend to colonise a larger volume of the visual feature space.
This pattern is consistent with the "niche filling" model \parencite{Simpson1944Tempo},
where speciation events are often associated with the exploration of novel regions
in the phenotypic landscape to minimise competition.
However, this expansion exhibits diminishing returns (\cref{fig:models_comparison}),
suggesting that the marginal gain in disparity decreases as taxa become extremely large.

In contrast, structural heterogeneity appears to be constrained that is largely independent of diversity.
At the order level, angle variance is statistically independent of clade size ($\rho=0.24,p>0.05$),
and at the family level, the correlation is weak ($\rho=0.34$).
This implies that even as a clade radiates swiftly, its internal geometric uniformity
does not change chaotically but remains confined within a stable range.

In the top-5 families with high disparity residuals, four of them are newly recognised based on
molecular phylogeny, indicating that they exhibit greater heterogeneity than long-established taxa,
which is why they remained unrecognised by zoologists for a long time.

Vangidae is a textbook example of the radiation evolution.
Although there aren't many species (39), they have evolved a variety of forms on Madagascar,
ranging from shrike-like and flycatcher-like to finch-like species \parencite{Reddy2012Diversification, Jonsson2012Ecological}.
The model captures this extreme morphological difference, producing a very high positive residual.

Similarly, Cotingidae is given the seventh-highest disparity residual among families.
Their differences in appearance are extremely eye-catching, ranging from the
brightly coloured cocks-of-the-rock (\textit{Rupicola})
to the relatively dull-coloured pihas (\textit{Lipaugus} and \textit{Snowornis}),
with huge differences \parencite{Winkler2020Cotingas}.

The disparity-through-time (DTT) analysis (\cref{fig:dtt}) exhibits a convex shape,
suggesting that morphological disparity increased rapidly after the K-Pg mass extinction.
This furtherly supports \posscite{Simpson1944Tempo} adaptive radiation and the niche filling hypothesis.
After mass extinctions vacated ecological niches,
birds rapidly evolved extreme morphological diversity to occupy these positions.
The surviving avian lineages rapidly radiated to fill the ecological niches,
establishing the major disparate body plans (visual morphotypes) within a short geological window.
The positive Morphological Disparity Index (MDI), evidenced in \cref{fig:dtt},
reinforces that the initial partitioning of the visual morphospace
was significantly faster than neutral drift expectations.